% =============================================================================
% Seção 4 — Avaliação
% =============================================================================

\section{Avaliação}
\label{sec:avaliacao}

Avaliamos o NeuroAssist AI por meio de um estudo de caso preliminar envolvendo um
paciente com ELA em estágio avançado e cinco cuidadores e voluntários. O estudo
mediu a latência ponta-a-ponta do sistema, a taxa de detecção de gesto deliberado e
a taxa de falsos positivos em condições naturalísticas. Os resultados devem ser
interpretados como um estudo de caso preliminar, e não como um ensaio clínico (ver
Seção~\ref{sec:discussao}).

\subsection{Participantes e Ética}
\label{sec:participantes}

Seis participantes tomaram parte no estudo (Tabela~\ref{tab:participantes}).
P0 é um paciente com ELA em estágio avançado com histórico clínico documentado;
P1--P3 são cuidadores familiares com contato regular e íntimo com P0; P4--P5 são
voluntários saudáveis sem contato prévio com o paciente, servindo como controles
de falsos positivos naturalísticos.

\begin{table}[htbp]
  \renewcommand{\arraystretch}{1.15}
  \caption{Participantes do Estudo}
  \label{tab:participantes}
  \centering
  \begin{tabular}{llll}
    \toprule
    ID & Papel & Relação com P0 & Sessão \\
    \midrule
    P0 & Paciente ELA          & ---                              & Calibração + detecção \\
    P1 & Cônjuge / cuidadora doméstica & Íntima, diária (fisioterapeuta)   & FP + alarme real \\
    P2 & Cuidador 1                   & Profissional NDIS, frequente      & FP + alarme real \\
    P3 & Cuidador 2                   & Profissional NDIS, frequente      & FP + alarme real \\
    P4 & Controle           & Nenhum                     & Baseline FP naturalístico \\
    P5 & Controle           & Nenhum                     & Baseline FP naturalístico \\
    \bottomrule
  \end{tabular}
\end{table}

Todos os participantes forneceram consentimento informado antes de sua sessão. O
estudo foi conduzido no domicílio de P0 em Melbourne, Austrália. P1 é cônjuge de P0
e fisioterapeuta, prestando cuidado domiciliar diário (não elegível para financiamento
NDIS como familiar, conforme política australiana); P2 e P3 são profissionais de suporte
financiados pelo NDIS (National Disability Insurance Scheme) com contato regular e
estruturado com P0.
O estudo constitui uma avaliação piloto; nenhuma aprovação formal de comitê de ética
em pesquisa (CEP) foi solicitada para este escopo de estudo de caso preliminar. A avaliação clínica completa com uma amostra
estatisticamente representativa de pacientes com ELA exigirá aprovação do CEP e está
planejada como trabalho futuro.

P0 apresenta ELA de início troncular (início dos sintomas: fim de 2019; duração da
doença na avaliação: aproximadamente 6,5~anos), fraqueza em quatro membros com
movimento voluntário palpebral preservado, e ventilação não invasiva (VNI) prescrita
desde o início de 2026 e ativa no momento do estudo. A capacidade vital forçada foi de
71\,\% e a pressão inspiratória nasal de esforço (SNIP) de 31\,\% na última avaliação
respiratória documentada antes do experimento (maio de 2023), com progressão lenta da
doença confirmada naquela ocasião. \fillme{Confirmar antes da submissão: escore
ALSFRS-R se disponível; avaliação atual da mobilidade palpebral.} Este perfil clínico
embasou o protocolo de calibração descrito na Seção~\ref{sec:calibracao}.

\subsection{Configuração Experimental}
\label{sec:configuracao}

\textbf{Hardware.} Cada sessão utilizou uma webcam USB comum
(\fillme{modelo/resolução, e.g., Logitech C920, 1080p@30fps}) posicionada na altura
dos olhos a aproximadamente 50--70\,cm do rosto do participante, conectada a um
notebook executando a pilha NeuroAssist AI localmente. Nenhum acesso externo à rede
foi necessário durante as sessões.

\textbf{Software.} O pipeline de detecção de borda foi executado como
\texttt{poc-edge-agent-webcam.py} (Python~3.11), comunicando-se com a aplicação
Angular (\texttt{poc-gesture-detector.ts}) via \texttt{localhost}. Os microsserviços
\texttt{evnt-service} e \texttt{gest-service} foram executados localmente via Docker
Compose. Scripts de medição registraram todos os timestamps em CSV; os timestamps de
entrega de alarme foram capturados por um hook Angular \texttt{effect()} que grava
em \texttt{localStorage} a cada renderização de alarme.

\textbf{Ambiente.} As sessões foram conduzidas no domicílio de P0 durante
\fillme{período do dia}. A iluminação ambiente era natural/artificial de interiores.

\subsection{Calibração do Paciente}
\label{sec:calibracao}

Antes dos ensaios formais, o \texttt{DetectionProfile} de P0 foi calibrado por meio
da observação do movimento palpebral natural por aproximadamente 5--10~minutos. O
perfil clínico de P0 — ELA progressiva com envolvimento respiratório documentado
(CVF 71\,\%, SNIP 31\,\%) e VNI prescrita desde o início de 2026 — é representativo
da população para a qual o NeuroAssist AI é primordialmente
projetado~\cite{berlowitz2021niv, yang2026pva}, e o DetectionProfile calibrado reflete
a mobilidade palpebral residual característica desse estágio. O procedimento de
calibração ajustou dois parâmetros: $\theta_\text{close}$ (limiar EAR para início da
piscada) e \texttt{BLINK\_MIN\_FRAMES} (quadros consecutivos mínimos abaixo de
$\theta_\text{close}$ necessários para registrar uma piscada).

A Tabela~\ref{tab:calibracao} reporta os parâmetros calibrados para P0 ao lado dos
padrões utilizados para P4 e P5.

\begin{table}[htbp]
  \renewcommand{\arraystretch}{1.15}
  \caption{Parâmetros de Calibração do DetectionProfile}
  \label{tab:calibracao}
  \centering
  \begin{tabular}{lccc}
    \toprule
    Participante & $\theta_\text{close}$ & $\theta_\text{open}$ & Quadros mín. \\
    \midrule
    P0 (paciente ELA) & \textbf{[X,XX]} & \textbf{[X,XX]} & \textbf{[N]} \\
    P4--P5 (controles, padrão) & 0,20 & 0,27 & 3 \\
    \bottomrule
  \end{tabular}
\end{table}

\subsection{Latência}
\label{sec:latencia}

A latência ponta-a-ponta foi decomposta em dois segmentos: (i)~gesto-para-HTTP,
medido como o tempo decorrido entre a máquina de estados entrar em \texttt{DONE} e a
resposta do \texttt{POST /api/v0/report-ingest} ser recebida pelo agente de borda; e
(ii)~evento-para-alarme-visível, medido como o tempo decorrido entre o timestamp
\texttt{createdAt} na resposta do \texttt{evnt-service} e o valor
\texttt{alarmShownAtMs} registrado em \texttt{localStorage}. A Tabela~\ref{tab:latencia}
reporta as estatísticas de latência em todos os ensaios.

\begin{table}[htbp]
  \renewcommand{\arraystretch}{1.15}
  \caption{Latência Ponta-a-Ponta (ms)}
  \label{tab:latencia}
  \centering
  \begin{tabular}{lcccc}
    \toprule
    Segmento & $N$ & Média & $\sigma$ & Mediana \\
    \midrule
    Gesto $\rightarrow$ HTTP  & [N] & [X] & [Y] & [Z] \\
    Evento $\rightarrow$ alarme & [N] & [X] & [Y] & [Z] \\
    \textbf{Total ponta-a-ponta} & [N] & \textbf{[X]} & [Y] & [Z] \\
    \bottomrule
  \end{tabular}
\end{table}

A latência ponta-a-ponta média foi de \textbf{[X]}\,ms ($\sigma$\,=\,[Y]\,ms, mediana
[Z]\,ms, $N$\,=\,[N] ensaios). O segmento evento-para-alarme representou a maior
parcela da latência total, consistente com o intervalo de polling de 5 segundos
utilizado na aplicação de cuidadores do PoC (contribuição máxima: 5.000\,ms; ver
Seção~\ref{sec:discussao} para a mitigação WebSocket planejada). O segmento
gesto-para-HTTP foi dominado pela sobrecarga HTTP local e permaneceu abaixo de
[X]\,ms em todos os ensaios.

\subsection{Taxa de Falsos Positivos}
\label{sec:taxafp}

A taxa de falsos positivos (TFP) foi medida como o número de disparos de alarme não
intencionais por minuto de atividade facial contínua, tanto para a máquina de estados
de gesto deliberado de seis estados (ME) quanto para um limiar EAR simples como
baseline. Cada sessão teve duração de 10~minutos de atividade facial naturalística
sem que o participante tentasse acionar o gesto. A Tabela~\ref{tab:fp} reporta a TFP
por participante e por grupo.

\begin{table}[htbp]
  \renewcommand{\arraystretch}{1.15}
  \caption{Taxa de Falsos Positivos (disparos/min)}
  \label{tab:fp}
  \centering
  \begin{tabular}{lccc}
    \toprule
    Participante & Grupo & Limiar simples & Máquina de estados \\
    \midrule
    P1 & Cuidador & [A] & [B] \\
    P2 & Cuidador & [A] & [B] \\
    P3 & Cuidador & [A] & [B] \\
    P4 & Controle & [A] & [B] \\
    P5 & Controle & [A] & [B] \\
    \midrule
    \textbf{Média cuidadores} & & [A] & \textbf{[B]} \\
    \textbf{Média controles}  & & [A] & \textbf{[B]} \\
    \bottomrule
  \end{tabular}
\end{table}

A máquina de estados de seis estados reduziu a TFP de [A]\,disparos/min (limiar
simples) para [B]\,disparos/min em média entre todos os participantes, uma redução
relativa de [C]\,\%. Os participantes cuidadores (P1--P3), que interagiram com P0 em
contextos emocionalmente naturais, exibiram TFP \textbf{[maior/menor/comparável]} em
relação aos participantes controle (P4--P5) em tarefas neutras.

\subsection{Taxa de Detecção (Recall)}
\label{sec:taxadeteccao}

P0 realizou $N$\,=\,[N] tentativas de dupla piscada deliberada com pausas de pelo
menos 30\,s entre tentativas, atingindo taxa de detecção de \textbf{[D]\,\%} ([M] de
[N] tentativas detectadas) com o \texttt{DetectionProfile} calibrado. Os participantes
controle P4 e P5 atingiram taxas de [D1]\,\% e [D2]\,\%, respectivamente, usando
parâmetros padrão. A taxa de detecção superior dos voluntários saudáveis, se
confirmada, é consistente com maior mobilidade palpebral residual em comparação com
a ELA em estágio avançado.
